% !Mode:: "TeX:UTF-8"
% !TEX root = ../../main.tex
% ============================================================
% TOPIC 7
% ============================================================
\mashgulot{7}{Cramer’s rule for systems of linear algebraic equations}

\begin{tarif}
Consider a system of $n$ linear equations in $n$ unknowns
\[\kf{\begin{cases}
a_{11}x_1+a_{12}x_2+\dots+a_{1n}x_n=b_1,\\
a_{21}x_1+a_{22}x_2+\dots+a_{2n}x_n=b_2,\\
\dots\dots\dots\dots\dots\dots\dots\dots\dots\dots\\
a_{n1}x_1+a_{n2}x_2+\dots+a_{nn}x_n=b_n.
\end{cases}}\]
The determinant $\Delta=\det A$ of its coefficient matrix is called the
\textbf{determinant of the system}. The determinant obtained from $\Delta$ by
replacing its $j$th column with the column of constant terms
$b=(b_1;\dots;b_n)^T$ is denoted by $\Delta_j$; we call it the $j$th
\textbf{auxiliary determinant}.
\end{tarif}

\begin{teorema}[Cramer’s rule]
If the determinant of the system is nonzero ($\Delta\neq0$), the system has a
unique solution, given by
\begin{equation}
x_j=\frac{\Delta_j}{\Delta},\qquad j=1,2,\dots,n.
\end{equation}
\end{teorema}

In practice one first computes $\Delta$. If $\Delta=0$, Cramer’s rule does not
apply: the system is then either inconsistent or has infinitely many solutions.
If $\Delta\neq0$, one computes the auxiliary determinants
$\Delta_1,\Delta_2,\dots,\Delta_n$, finds the unknowns from the formulas
$x_j=\Delta_j/\Delta$, and checks the result by substituting it into the system.

\subsection{Cramer’s rule for $2\times2$ systems}

For the system
$\begin{cases}a_{11}x_1+a_{12}x_2=b_1,\\ a_{21}x_1+a_{22}x_2=b_2\end{cases}$
the three determinants are
\[\kf{\Delta=\begin{vmatrix}a_{11}&a_{12}\\a_{21}&a_{22}\end{vmatrix},\quad
\Delta_1=\begin{vmatrix}b_1&a_{12}\\b_2&a_{22}\end{vmatrix},\quad
\Delta_2=\begin{vmatrix}a_{11}&b_1\\a_{21}&b_2\end{vmatrix}.}\]

\begin{misol}
Solve the system $\begin{cases}3x_1+2x_2=7,\\ x_1-x_2=-1\end{cases}$ by Cramer’s rule.

\textbf{Solution.}
\[\kf{\Delta=\begin{vmatrix}3&2\\1&-1\end{vmatrix}=-3-2=-5\neq0,\quad
\Delta_1=\begin{vmatrix}7&2\\-1&-1\end{vmatrix}=-7+2=-5,\quad
\Delta_2=\begin{vmatrix}3&7\\1&-1\end{vmatrix}=-3-7=-10.}\]
Hence $x_1=\dfrac{-5}{-5}=1$ and $x_2=\dfrac{-10}{-5}=2$.

\textbf{Check:} $3\cdot1+2\cdot2=7\;\checkmark$; \ $1-2=-1\;\checkmark$.
\end{misol}

\begin{misol}
Solve the system $\begin{cases}2x_1-5x_2=4,\\ 3x_1+x_2=-2.\end{cases}$

\textbf{Solution.} We have $\Delta=2\cdot1-(-5)\cdot3=17$, \
$\Delta_1=4\cdot1-(-5)(-2)=-6$ and $\Delta_2=2\cdot(-2)-4\cdot3=-16$, so
\[\kf{x_1=-\frac{6}{17},\qquad x_2=-\frac{16}{17}.}\]
\end{misol}

\subsection{Cramer’s rule for $3\times3$ systems}

\begin{misol}
Solve the system
$\begin{cases}x_1+x_2+x_3=6,\\ 2x_1-x_2+x_3=3,\\ x_1+2x_2-x_3=2\end{cases}$
by Cramer’s rule.

\textbf{Solution.} Expanding the determinant of the system along the first row,
we find
\[\kf{\Delta=\begin{vmatrix}1&1&1\\2&-1&1\\1&2&-1\end{vmatrix}
=1(1-2)-1(-2-1)+1(4+1)=-1+3+5=7\neq0.}\]
The auxiliary determinants are
\[\kf{\Delta_1=\begin{vmatrix}6&1&1\\3&-1&1\\2&2&-1\end{vmatrix}=7,\qquad
\Delta_2=\begin{vmatrix}1&6&1\\2&3&1\\1&2&-1\end{vmatrix}=14,\qquad
\Delta_3=\begin{vmatrix}1&1&6\\2&-1&3\\1&2&2\end{vmatrix}=21.}\]
Therefore $x_1=\dfrac{7}{7}=1$, $x_2=\dfrac{14}{7}=2$, $x_3=\dfrac{21}{7}=3$.

\textbf{Check:} $1+2+3=6$, \ $2-2+3=3$, \ $1+4-3=2$ \ $\checkmark$.
\end{misol}

\subsection{The case $\Delta=0$}

\begin{teorema}
If $\Delta=0$ and at least one of the auxiliary determinants is nonzero, the
system is inconsistent. If $\Delta=0$ and all $\Delta_j=0$, the system is either
inconsistent or has infinitely many solutions; in this case it has to be
investigated by Gaussian elimination.
\end{teorema}

\begin{misol}
For the system $\begin{cases}x_1+2x_2=3,\\ 2x_1+4x_2=7\end{cases}$ we have
$\Delta=4-4=0$ and $\Delta_1=3\cdot4-2\cdot7=-2\neq0$. The system is therefore
\textbf{inconsistent}: its equations contradict each other.
\end{misol}

\begin{misol}
For the system $\begin{cases}x_1-3x_2=2,\\ -2x_1+6x_2=-4\end{cases}$ we have
$\Delta=6-6=0$, $\Delta_1=12-12=0$ and $\Delta_2=-4+4=0$. The second equation is
the first one multiplied by $-2$, so the system has \textbf{infinitely many
solutions}: $x_1=2+3t$, $x_2=t$, $t\in\R$.
\end{misol}

\begin{eslatma}
For $n$ unknowns Cramer’s rule requires $n+1$ determinants of order $n$. It is
therefore convenient mainly for $n=2$ or $3$ and in theoretical work, for
example when a system depends on parameters. For large systems Gaussian
elimination is far cheaper.
\end{eslatma}

\subsection{Applications of Cramer’s rule in information security}

\textbf{Recovering a block of the Hill cipher.}
In the Hill cipher a plaintext block $p$ is encrypted as $c=Kp$. If the key
matrix $K$ is known, recovering $p$ amounts to solving the linear system $Kp=c$.
For small keys ($2\times2$ or $3\times3$) this can be done directly by Cramer’s
rule, and the condition $\Delta=\det K\neq0$ is a necessary condition for the
key to be admissible.

\begin{misol}
Suppose that a block encrypted with the key $K=\begin{pmatrix}3&2\\5&7\end{pmatrix}$
(without reduction modulo $26$) is $c=\begin{pmatrix}29\\63\end{pmatrix}$. We solve
$Kp=c$:
\[\kf{\Delta=21-10=11,\quad
\Delta_1=\begin{vmatrix}29&2\\63&7\end{vmatrix}=203-126=77,\quad
\Delta_2=\begin{vmatrix}3&29\\5&63\end{vmatrix}=189-145=44.}\]
Thus $p_1=77/11=7$ and $p_2=44/11=4$, i.e. $p=(7;4)$, which corresponds to the
letters “H” and “E” ($A=0,\dots,Z=25$).
\end{misol}

\textbf{Key recovery in a known-plaintext attack.}
If an attacker knows several plaintext–ciphertext pairs, the unknown entries of
the key satisfy a system of linear equations. When the determinant of this
system is nonzero, Cramer’s rule recovers the key uniquely. Because the Hill
cipher is \textbf{linear}, it offers no resistance to a known-plaintext attack.

\textbf{Systems with parameters and the key space.}
For a family of keys depending on a parameter, the roots of the equation
$\Delta(k)=0$ give the inadmissible keys. For instance, for
$K=\begin{pmatrix}1&k\\k&4\end{pmatrix}$ we get $\Delta=4-k^2$, so the values
$k=\pm2$ cannot be used.

\subsection{Computing with Python, Maple and Mathcad}
We solve the $3\times3$ system of the third example once more. In Python it is
enough to replace the columns of $A$ by $b$ in turn and divide the determinants:
\begin{Verbatim}[fontsize=\footnotesize,frame=single,label={\small Python}]
from sympy import Matrix
A = Matrix([[1, 1, 1], [2, -1, 1], [1, 2, -1]])
b = Matrix([6, 3, 2])
D = A.det()
x = []
for j in range(3):
    Aj = A.copy()
    Aj[:, j] = b              # j-th column replaced by b
    x.append(Aj.det() / D)
print(D, x)
\end{Verbatim}
\begin{Verbatim}[fontsize=\footnotesize,frame=single,label={\small Output}]
7 [1, 2, 3]
\end{Verbatim}
Maple does the same with \texttt{Determinant} and \texttt{Column}; the command
\texttt{LinearSolve} serves as a check:
\begin{Verbatim}[fontsize=\footnotesize,frame=single,label={\small Maple}]
with(LinearAlgebra):
A := Matrix([[1, 1, 1], [2, -1, 1], [1, 2, -1]]):
b := Vector([6, 3, 2]):
D0 := Determinant(A);
D1 := Determinant(<b | Column(A, 2) | Column(A, 3)>);
D2 := Determinant(<Column(A, 1) | b | Column(A, 3)>);
D3 := Determinant(<Column(A, 1) | Column(A, 2) | b>);
D1/D0, D2/D0, D3/D0;
LinearSolve(A, b);
\end{Verbatim}
In Mathcad Prime the auxiliary matrices are built from the columns of $A$
(the operator $A^{\langle j\rangle}$, typed with Ctrl+Shift+C) and the function
\texttt{augment}; $|A|$ is the determinant, and \texttt{lsolve(A, b)} gives the
solution directly (with \texttt{ORIGIN = 0} the columns are numbered from zero):
\begin{Verbatim}[fontsize=\footnotesize,frame=single,label={\small Mathcad}]
     | 1   1   1 |         | 6 |
A := | 2  -1   1 |    b := | 3 |
     | 1   2  -1 |         | 2 |
A1 := augment(b, A<1>, A<2>)
A2 := augment(A<0>, b, A<2>)
A3 := augment(A<0>, A<1>, b)
|A| = 7   |A1| = 7   |A2| = 14   |A3| = 21
x1 := |A1|/|A|   x2 := |A2|/|A|   x3 := |A3|/|A|
x1 = 1           x2 = 2           x3 = 3
                 | 1 |
lsolve(A, b) =   | 2 |
                 | 3 |
\end{Verbatim}

\subsection{Exercises on Cramer’s rule}

\emph{1. Solve the following $2\times 2$ systems by Cramer’s rule.}
\bigskip
\begin{multicols}{2}\raggedright
\begin{enumerate}[label=\arabic*., leftmargin=*]
\item $\begin{cases}-x_1-4x_2=26\\4x_2=-20\end{cases}$
\item $\begin{cases}2x_1-5x_2=-22\\-x_1+3x_2=12\end{cases}$
\item $\begin{cases}2x_1=-10\\-6x_1+3x_2=21\end{cases}$
\item $\begin{cases}-2x_1=10\\-4x_1+2x_2=26\end{cases}$
\item $\begin{cases}-3x_1-x_2=-10\\-5x_1+2x_2=-35\end{cases}$
\item $\begin{cases}3x_1-6x_2=-21\\3x_1-3x_2=-9\end{cases}$
\item $\begin{cases}2x_1=2\\6x_1-x_2=3\end{cases}$
\item $\begin{cases}x_1-x_2=10\\-2x_1-3x_2=0\end{cases}$
\item $\begin{cases}5x_1+6x_2=3\\-3x_1-5x_2=1\end{cases}$
\item $\begin{cases}2x_1+x_2=0\\-x_1+5x_2=-11\end{cases}$
\item $\begin{cases}3x_1+6x_2=9\\-x_1-x_2=-4\end{cases}$
\item $\begin{cases}4x_1+3x_2=7\\4x_1+x_2=-3\end{cases}$
\item $\begin{cases}4x_2=-4\\-x_1-6x_2=5\end{cases}$
\item $\begin{cases}-4x_1+3x_2=15\\-5x_1+x_2=27\end{cases}$
\item $\begin{cases}-3x_1=15\\x_2=-4\end{cases}$
\item $\begin{cases}x_1=-4\\2x_1-2x_2=-8\end{cases}$
\item $\begin{cases}2x_1-2x_2=-10\\5x_1=-5\end{cases}$
\item $\begin{cases}-3x_2=12\\-4x_1-5x_2=36\end{cases}$
\item $\begin{cases}x_1+3x_2=-20\\-4x_1-2x_2=20\end{cases}$
\item $\begin{cases}-4x_1=-12\\2x_1-x_2=3\end{cases}$
\item $\begin{cases}3x_1+4x_2=-14\\4x_1+5x_2=-19\end{cases}$
\item $\begin{cases}6x_1+4x_2=0\\6x_1+2x_2=0\end{cases}$
\item $\begin{cases}4x_1=-20\\-6x_1-3x_2=39\end{cases}$
\item $\begin{cases}3x_1-4x_2=-15\\2x_1-5x_2=-17\end{cases}$
\item $\begin{cases}-x_1+3x_2=-10\\-x_1+x_2=0\end{cases}$
\item $\begin{cases}-2x_1-5x_2=-5\\-4x_1-5x_2=-15\end{cases}$
\item $\begin{cases}2x_1-x_2=10\\-4x_1+5x_2=-38\end{cases}$
\item $\begin{cases}-x_1-4x_2=-5\\-x_1+6x_2=15\end{cases}$
\item $\begin{cases}-3x_1-3x_2=-12\\2x_1+x_2=3\end{cases}$
\item $\begin{cases}x_1-2x_2=5\\-3x_1+5x_2=-14\end{cases}$
\end{enumerate}
\end{multicols}

\emph{2. Solve the following $3\times 3$ systems by Cramer’s rule.}
\bigskip
\begin{multicols}{2}\raggedright
\begin{enumerate}[label=\arabic*., leftmargin=*]
\item $\begin{cases}3x_1+x_2+x_3=5\\-3x_1-x_2-3x_3=-11\\-x_1+3x_2-x_3=-7\end{cases}$
\item $\begin{cases}-4x_1+3x_2+x_3=-8\\-3x_1-3x_2+2x_3=-15\\-x_1+3x_2-2x_3=7\end{cases}$
\item $\begin{cases}2x_1+3x_2+2x_3=6\\-3x_1-2x_2-2x_3=-11\\-2x_1-4x_2-2x_3=-4\end{cases}$
\item $\begin{cases}-x_1+x_2=-2\\4x_1+2x_2-2x_3=12\\-4x_1+x_2+3x_3=-2\end{cases}$
\item $\begin{cases}4x_1+4x_2+3x_3=-16\\-3x_1+4x_2-4x_3=-16\\-x_1-x_2=7\end{cases}$
\item $\begin{cases}4x_1-x_2+3x_3=-5\\-2x_1+2x_2-3x_3=-2\\2x_1+3x_2+x_3=-7\end{cases}$
\item $\begin{cases}-3x_1-x_2=-8\\-3x_1-2x_2+x_3=-10\\-2x_1-2x_3=0\end{cases}$
\item $\begin{cases}2x_1+3x_2-2x_3=-1\\-x_1-2x_2+2x_3=2\\4x_1+2x_2+x_3=7\end{cases}$
\item $\begin{cases}x_1-3x_2+x_3=3\\-4x_1+x_2+4x_3=9\\3x_1+3x_2-4x_3=-7\end{cases}$
\item $\begin{cases}4x_1+3x_2+x_3=-15\\-3x_1-4x_3=12\\-2x_1+2x_2-3x_3=1\end{cases}$
\item $\begin{cases}-3x_2-x_3=-6\\-3x_1-3x_3=-12\\3x_1-4x_2+x_3=4\end{cases}$
\item $\begin{cases}-2x_1-4x_2+4x_3=2\\-x_1-3x_2-2x_3=1\\-4x_2-2x_3=0\end{cases}$
\item $\begin{cases}-4x_2-4x_3=0\\-4x_1+4x_2+4x_3=4\\-x_1+4x_2+3x_3=4\end{cases}$
\item $\begin{cases}2x_1+3x_2+4x_3=0\\2x_1+4x_2=-10\\-x_1-x_2+x_3=5\end{cases}$
\item $\begin{cases}-3x_1+2x_2+4x_3=18\\-x_2=0\\-4x_1+3x_2-2x_3=2\end{cases}$
\item $\begin{cases}-4x_1+x_3=2\\x_1+4x_2+x_3=1\\-x_1-4x_2=-3\end{cases}$
\item $\begin{cases}-4x_1+x_2+2x_3=-26\\-3x_1+3x_2=-24\\4x_1-x_2-x_3=23\end{cases}$
\item $\begin{cases}-x_1-3x_2+4x_3=30\\-2x_1+2x_2+x_3=0\\3x_1-2x_2=2\end{cases}$
\item $\begin{cases}2x_1+4x_2-2x_3=-18\\4x_1+4x_2-4x_3=-28\\-x_1-3x_2-4x_3=6\end{cases}$
\item $\begin{cases}-3x_1+2x_2+3x_3=1\\4x_1-4x_2-4x_3=4\\4x_1-x_2+3x_3=13\end{cases}$
\item $\begin{cases}-x_2-x_3=5\\-x_1+3x_2+3x_3=-15\\2x_1-3x_2+3x_3=9\end{cases}$
\item $\begin{cases}-3x_1-2x_2+x_3=-12\\-2x_3=6\\-4x_1+3x_2-4x_3=0\end{cases}$
\item $\begin{cases}4x_1-x_2=13\\-3x_1+3x_2-4x_3=-3\\3x_2-3x_3=9\end{cases}$
\item $\begin{cases}2x_1-x_2-x_3=-8\\-3x_1-3x_2-2x_3=-6\\4x_1+x_3=-8\end{cases}$
\item $\begin{cases}-3x_1+x_2-x_3=-8\\3x_1+3x_2+2x_3=22\\-4x_1-2x_2-4x_3=-26\end{cases}$
\item $\begin{cases}-2x_2+2x_3=2\\x_1+2x_2+x_3=5\\-3x_1+x_2-4x_3=-10\end{cases}$
\item $\begin{cases}-3x_1-x_2-4x_3=0\\x_3=2\\-3x_1+2x_2+2x_3=15\end{cases}$
\item $\begin{cases}2x_1+4x_2+x_3=1\\-x_1+x_2+2x_3=-11\\-4x_1+2x_2+4x_3=-30\end{cases}$
\item $\begin{cases}-4x_1+2x_2+3x_3=20\\-2x_1+3x_3=10\\-4x_1+4x_2-2x_3=16\end{cases}$
\item $\begin{cases}4x_1-x_2+3x_3=-7\\x_1+3x_2+2x_3=-14\\-2x_1+4x_2-x_3=-8\end{cases}$
\end{enumerate}
\end{multicols}

\emph{3. Applied problems (10 problems).}
\bigskip
\begin{enumerate}[label=\arabic*., leftmargin=*]
\item \textbf{(Hill cipher.)} A plaintext block $p$ was encrypted as $c=Kp$ (without reduction modulo $26$) with the key matrix $K=\begin{pmatrix}3 & 2\\ 5 & 7\end{pmatrix}$. Given the ciphertext $c=\begin{pmatrix}29\\ 63\end{pmatrix}$, solve the system $Kp=c$ by Cramer’s rule and find the plaintext $p$.
\item \textbf{(Hill cipher.)} A plaintext block $p$ was encrypted as $c=Kp$ (without reduction modulo $26$) with the key matrix $K=\begin{pmatrix}2 & 3\\ 1 & 4\end{pmatrix}$. Given the ciphertext $c=\begin{pmatrix}37\\ 31\end{pmatrix}$, solve the system $Kp=c$ by Cramer’s rule and find the plaintext $p$.
\item \textbf{(Hill cipher.)} A plaintext block $p$ was encrypted as $c=Kp$ (without reduction modulo $26$) with the key matrix $K=\begin{pmatrix}5 & 1\\ 3 & 2\end{pmatrix}$. Given the ciphertext $c=\begin{pmatrix}59\\ 62\end{pmatrix}$, solve the system $Kp=c$ by Cramer’s rule and find the plaintext $p$.
\item \textbf{(A parameter.)} For which values of $k$ does the system $\begin{cases}x_1+kx_2=3\\kx_1+4x_2=6\end{cases}$ have a unique solution? Find the solution for $k=1$.
\item \textbf{(A parameter.)} For which values of $a$ can the system $\begin{cases}ax_1+2x_2=1\\3x_1+(a+1)x_2=2\end{cases}$ not be solved by Cramer’s rule?
\item \textbf{(Determinants.)} For the system $\begin{cases}2x_1-x_2+x_3=5\\x_1+3x_2-2x_3=-6\\3x_1+x_2+4x_3=10\end{cases}$ compute the determinants $\Delta,\Delta_1,\Delta_2,\Delta_3$ and find the solution.
\item \textbf{(Network traffic.)} The numbers $x_1,x_2,x_3$ of packets (in thousands) that passed through three servers satisfy the system $\begin{cases}x_1+x_2+x_3=12\\2x_1-x_2+x_3=9\\x_1+2x_2-x_3=5\end{cases}$. Find them by Cramer’s rule.
\item \textbf{(A fractional solution.)} Solve the system $\begin{cases}3x_1+2x_2=7\\4x_1-5x_2=3\end{cases}$ by Cramer’s rule.
\item \textbf{(The case $\Delta=0$.)} Investigate the systems $\begin{cases}2x_1-4x_2=6\\-x_1+2x_2=-3\end{cases}$ and $\begin{cases}2x_1-4x_2=6\\-x_1+2x_2=1\end{cases}$.
\item \textbf{(Key recovery.)} The key $K=\begin{pmatrix}a&b\\c&d\end{pmatrix}$ of a Hill cipher is unknown, but it is known that $K\binom{1}{0}=\binom{3}{5}$ and $K\binom{1}{1}=\binom{5}{12}$. Find $a,b,c,d$ using Cramer’s rule.
\end{enumerate}

